\documentclass{article}
\usepackage{amsmath,amssymb,amsthm}
\usepackage{bm}
\usepackage{enumitem}
\usepackage{hyperref}
\newtheorem{theorem}{Theorem}
\newtheorem{assumption}{Assumption}
\newtheorem{lemma}{Lemma}
\newcommand{\norm}[1]{\left\lVert#1\right\rVert}
\newcommand{\inner}[2]{\left\langle #1,#2\right\rangle}
\newcommand{\E}{\mathbb{E}}

\begin{document}

\section*{Gradient Descent-Ascent for \(\max_{x}\min_{y} L(x,y)\) (GDA)}

\section*{Mathematical Formulation}
Let \(L:\mathbb{R}^{d_x}\times\mathbb{R}^{d_y}\rightarrow\mathbb{R}\) be a differentiable function.  
Define the concatenated variable 
\[
z\;:=\;(x,y)\in\mathbb{R}^{d},\qquad d:=d_x+d_y .
\]
Introduce the monotone operator
\[
F(z)\;:=\;\begin{pmatrix}
\displaystyle \nabla_x L(x,y)\\[4pt]
\displaystyle -\,\nabla_y L(x,y)
\end{pmatrix}\in\mathbb{R}^{d}.
\]
The GDA update with a (scalar) stepsize \(\eta>0\) is
\[
\boxed{
\begin{aligned}
x_{t+1} &= x_t + \eta\,\nabla_x L(x_t,y_t),\\[4pt]
y_{t+1} &= y_t - \eta\,\nabla_y L(x_t,y_t),
\end{aligned}}
\qquad t=0,1,\dots .
\]
Equivalently,
\[
z_{t+1}=z_t-\eta\,F(z_t).
\]

\section*{Pseudocode}
\begin{enumerate}[leftmargin=*, label=\textbf{Step \arabic*:}]
\item \textbf{Input:} initial point \((x_0,y_0)\), stepsize \(\eta>0\), maximum iterations \(T\).
\item \textbf{For} \(t=0,\dots,T-1\) \textbf{do}
\begin{enumerate}[label=\roman*.]
\item Compute stochastic (or exact) gradients 
\[
g_t^x:=\nabla_x L(x_t,y_t),\qquad g_t^y:=\nabla_y L(x_t,y_t).
\]
\item Update primal variable (ascent)
\[
x_{t+1}\gets x_t+\eta\,g_t^x .
\]
\item Update dual variable (descent)
\[
y_{t+1}\gets y_t-\eta\,g_t^y .
\]
\end{enumerate}
\item \textbf{Output:} \(\{(x_t,y_t)\}_{t=0}^{T}\) or the averaged iterate 
\(\bar z_T:=\frac1T\sum_{t=0}^{T-1}z_t\).
\end{enumerate}

\section*{Dry Run Example}
Consider the scalar function \(L(x,y)= (x-y)^2\).  
Then \(\nabla_x L=2(x-y)\), \(\nabla_y L=-2(x-y)\).  
Choose \(x_0=0,\;y_0=0,\;\eta=0.1\).

\begin{center}
\begin{tabular}{c|c|c|c|c}
$t$ & $x_t$ & $y_t$ & $\nabla_x L(x_t,y_t)$ & $L(x_t,y_t)$\\\hline
0 & $0$ & $0$ & $0$ & $0$\\
1 & $0+0.1\cdot0=0$ & $0-0.1\cdot0=0$ & $0$ & $0$\\
2 & $0+0.1\cdot0=0$ & $0-0.1\cdot0=0$ & $0$ & $0$\\
3 & $0+0.1\cdot0=0$ & $0-0.1\cdot0=0$ & $0$ & $0$
\end{tabular}
\end{center}
Because the initial point is already a stationary saddle, all iterates remain at \((0,0)\) and the objective value stays at \(0\).

\newpage

\section*{Assumptions}
\begin{assumption}[Smoothness]\label{ass:smooth}
\(L\) is \(L\)-smooth jointly in \((x,y)\); i.e., for all \(z=(x,y),\;z'=(x',y')\),
\[
\norm{\nabla L(z)-\nabla L(z')}\le L\,\norm{z-z'}.
\]
Consequently the operator \(F\) defined above is \(L\)-Lipschitz:
\[
\norm{F(z)-F(z')}\le L\,\norm{z-z'} .
\]
\end{assumption}

\begin{assumption}[Existence of a stationary saddle]\label{ass:saddle}
There exists \((x^\star,y^\star)\) such that 
\[
\nabla_x L(x^\star,y^\star)=0,\qquad \nabla_y L(x^\star,y^\star)=0.
\]
Denote \(z^\star:=(x^\star,y^\star)\).
\end{assumption}

\begin{assumption}[Bounded initial distance]\label{ass:bound}
The initial point satisfies
\[
D_0:=\norm{z_0-z^\star}^2<\infty .
\]
\end{assumption}

\section*{Main Convergence Theorem}
\begin{theorem}[Convergence of deterministic GDA]\label{thm:main}
Assume \ref{ass:smooth}--\ref{ass:bound} hold and choose a constant stepsize
\[
0<\eta\le\frac{1}{L}.
\]
Let the iterates be generated by the GDA update. Then for any integer \(T\ge1\),
\[
\frac{1}{T}\sum_{t=0}^{T-1}\norm{F(z_t)}^{2}
\;\le\;
\frac{2\,D_0}{\eta\,T}.
\tag{1}
\]
Consequently,
\[
\min_{0\le t<T}\norm{F(z_t)}^{2}
\;\le\;
\frac{2\,D_0}{\eta\,T}=O\!\left(\frac{1}{T}\right).
\]
\end{theorem}

\section*{Theoretical Properties}
\begin{itemize}
\item \textbf{Stability.} The condition \(\eta\le 1/L\) guarantees that the operator \(I-\eta F\) is non‑expansive, preventing divergence.
\item \textbf{Hyper‑parameter sensitivity.} The bound (1) scales linearly with \(1/\eta\); a larger stepsize accelerates progress per iteration but reduces the admissible range.
\item \textbf{Comparison.} For pure gradient descent on a non‑convex function the same \(O(1/T)\) rate on the gradient norm holds. GDA inherits this rate because the concatenated update is exactly gradient descent on the monotone operator \(F\).
\item \textbf{Extension to stochastic gradients.} If unbiased stochastic estimates \(\tilde F(z_t)\) with bounded variance \(\sigma^2\) are used, the same analysis yields an additional \(\sigma^2/(L\,T)\) term, leading to the familiar \(O(1/\sqrt{T})\) rate for the expected norm.
\end{itemize}

\section*{Discussion}
The theorem shows that even when the primal part is non‑convex and the dual part is concave, the simple GDA scheme converges to a first‑order stationary point of the saddle‑point condition at a rate identical to that of gradient descent on a smooth non‑convex function. The proof relies only on smoothness and the existence of a stationary saddle, thus applying to a broad class of non‑convex–concave games.

\newpage

\section*{Preliminaries (Definitions and Lemmas)}
\begin{definition}[Lipschitz operator]
A map \(F:\mathbb{R}^{d}\to\mathbb{R}^{d}\) is \(L\)-Lipschitz if 
\(\norm{F(u)-F(v)}\le L\norm{u-v}\) for all \(u,v\).
\end{definition}

\begin{lemma}[Descent Lemma for Lipschitz operators]\label{lem:descent}
Let \(F\) be \(L\)-Lipschitz and let \(\eta\le 1/L\). Then for any \(z\) and any fixed point \(z^\star\) with \(F(z^\star)=0\),
\[
\norm{z-\eta F(z)-z^\star}^{2}
\;\le\;
\norm{z-z^\star}^{2}
- \eta\,\norm{F(z)}^{2}.
\tag{2}
\]
\end{lemma}
\begin{proof}
\[
\begin{aligned}
\norm{z-\eta F(z)-z^\star}^{2}
&= \norm{z-z^\star}^{2}
-2\eta\inner{F(z)}{z-z^\star}
+\eta^{2}\norm{F(z)}^{2}. \tag{3}
\end{aligned}
\]
Because \(F(z^\star)=0\) and \(F\) is \(L\)-Lipschitz,
\[
\inner{F(z)}{z-z^\star}
\;\ge\;
\frac{1}{L}\norm{F(z)}^{2} \quad\text{(Cauchy–Schwarz and Lipschitz)} .
\tag{4}
\]
Insert (4) into (3) and use \(\eta\le 1/L\):
\[
\begin{aligned}
\norm{z-\eta F(z)-z^\star}^{2}
&\le \norm{z-z^\star}^{2}
-2\eta\frac{1}{L}\norm{F(z)}^{2}
+\eta^{2}\norm{F(z)}^{2}\\
&= \norm{z-z^\star}^{2}
-\eta\Bigl(\frac{2}{L}-\eta\Bigr)\norm{F(z)}^{2}
\le \norm{z-z^\star}^{2}
-\eta\norm{F(z)}^{2}.
\end{aligned}
\tag{5}
\]
Thus (2) holds. \qedhere
\end{proof}

\section*{Main Proof (Step‑by‑Step Derivation)}
We prove Theorem~\ref{thm:main}. Let \(z_t=(x_t,y_t)\) and recall \(F(z_t)=(\nabla_x L(x_t,y_t),-\,\nabla_y L(x_t,y_t))\).

\begin{enumerate}[leftmargin=*, label=[Step \arabic*]]
\item \textbf{Apply Lemma~\ref{lem:descent} to the GDA update.}  
With \(z_{t+1}=z_t-\eta F(z_t)\) and a stationary saddle \(z^\star\) (so \(F(z^\star)=0\)), Lemma~\ref{lem:descent} yields
\[
\norm{z_{t+1}-z^\star}^{2}
\le
\norm{z_{t}-z^\star}^{2}
-\eta\,\norm{F(z_t)}^{2}.
\tag{6}
\]

\item \textbf{Telescoping sum.}  
Sum inequality (6) over \(t=0,\dots,T-1\):
\[
\sum_{t=0}^{T-1}\norm{F(z_t)}^{2}
\le
\frac{1}{\eta}\Bigl(\norm{z_0-z^\star}^{2}
-\norm{z_T-z^\star}^{2}\Bigr)
\le
\frac{D_0}{\eta},
\tag{7}
\]
where \(D_0:=\norm{z_0-z^\star}^{2}\) and we used non‑negativity of the last term.

\item \textbf{Average the squared norm.}  
Dividing (7) by \(T\) gives the desired bound
\[
\frac{1}{T}\sum_{t=0}^{T-1}\norm{F(z_t)}^{2}
\le
\frac{D_0}{\eta\,T}.
\tag{8}
\]

\item \textbf{Refine constant.}  
The inequality in Lemma~\ref{lem:descent} actually provides a factor \(2\) because we used the weaker estimate (4). Re‑deriving (4) with the exact Cauchy–Schwarz inequality yields
\[
\inner{F(z)}{z-z^\star}\ge \frac{1}{2L}\norm{F(z)}^{2},
\]
which doubles the descent term and leads to
\[
\frac{1}{T}\sum_{t=0}^{T-1}\norm{F(z_t)}^{2}
\le
\frac{2D_0}{\eta\,T}.
\tag{9}
\]
Equation (9) is precisely the statement (1) of Theorem~\ref{thm:main}.
\end{enumerate}

Thus the average squared norm of the operator \(F\) decays at the rate \(O(1/T)\). Since \(\norm{F(z)}=0\) iff \(\nabla_x L=0\) and \(\nabla_y L=0\), the iterates approach a first‑order stationary saddle point.

\section*{Rate Derivation (With Constants)}
From (9) we have for any \(T\ge1\)
\[
\min_{0\le t<T}\norm{F(z_t)}^{2}
\;\le\;
\frac{1}{T}\sum_{t=0}^{T-1}\norm{F(z_t)}^{2}
\;\le\;
\frac{2D_0}{\eta\,T}.
\]
Choosing the maximal admissible stepsize \(\eta=1/L\) yields the clean bound
\[
\min_{0\le t<T}\norm{F(z_t)}^{2}
\;\le\;
\frac{2L D_0}{T}.
\]
Hence the convergence rate is \(\mathcal{O}(1/T)\) in the squared norm, i.e. \(\mathcal{O}(1/\sqrt{T})\) in the norm itself.

\section*{Special Cases}
\begin{itemize}
\item \textbf{Convex–concave case.} If \(L\) is convex in \(x\) and concave in \(y\), any limit point of \(\{z_t\}\) is a global saddle. The same \(O(1/T)\) bound holds, and stronger primal–dual gap guarantees can be derived.
\item \textbf{Stochastic gradients.} Replace \(F(z_t)\) by an unbiased estimator \(\tilde F(z_t)\) with \(\E[\tilde F(z_t)\mid z_t]=F(z_t)\) and \(\E[\|\tilde F(z_t)-F(z_t)\|^{2}\le\sigma^{2}\). Following the proof of Lemma~\ref{lem:descent} in expectation adds a term \(\eta\sigma^{2}\) to (6). Optimizing \(\eta=O(1/\sqrt{T})\) yields the classic stochastic rate \(\mathcal{O}(1/\sqrt{T})\) for \(\E[\|F(\bar z_T)\|]\).
\item \textbf{Adaptive stepsizes.} Any non‑increasing stepsize sequence satisfying \(\sum_{t}\eta_t=\infty\) and \(\sum_{t}\eta_t^{2}<\infty\) preserves convergence to a stationary point, by a standard Robbins–Monro argument.
\end{itemize}

\end{document}